% for_astronaut — user manual (LaTeX version)
%
% LaTeX edition of docs/manual.md, with the same content. It documents where the
% PDF manual comes from and lets anyone rebuild a PDF with a standard TeX
% distribution:
%
%     latexmk -pdf -outdir=build docs/manual.tex
%
% Written with AI assistance (Claude Code); see docs/ai_usage.md.
% Keep it in sync with docs/manual.md whenever the manual changes.

\documentclass[11pt,a4paper]{article}

\usepackage[T1]{fontenc}
\usepackage[utf8]{inputenc}
\usepackage{lmodern}
\renewcommand{\familydefault}{\sfdefault}  % sans-serif body, like the game UI
\usepackage[margin=2.2cm]{geometry}
\usepackage[table]{xcolor}
\usepackage{array}
\usepackage{tabularx}
\usepackage{booktabs}
\usepackage{enumitem}
\usepackage{fancyvrb}
\usepackage{fancyhdr}
\usepackage{graphicx}
\usepackage{tikz}
\usepackage[explicit]{titlesec}
\usepackage[hidelinks]{hyperref}

% The images are looked up from the repository root or from docs/.
\graphicspath{{./}{../}}

\hypersetup{
  pdftitle    = {for_astronaut --- user manual},
  pdfauthor   = {for_code UFRJ},
  pdfsubject  = {Written with AI assistance (Claude Code) and reviewed by the team. See docs/ai_usage.md.},
  pdfkeywords = {AI-assisted, Claude Code, NASA Space Apps 2026},
}

% Palette of the game art: navy background, violet buttons, yellow "_" in the logo.
\definecolor{navy}{HTML}{1B1441}
\definecolor{violet}{HTML}{614AD3}
\definecolor{logoyellow}{HTML}{FFFF00}
\definecolor{stargrey}{HTML}{322A3C}
\definecolor{softgrey}{HTML}{B9B4D6}
\colorlet{accent}{violet}

% Sections look like the violet menu buttons, with a yellow edge.
\titleformat{\section}
  {\normalfont\large\bfseries\color{white}}
  {}{0pt}
  {\colorbox{violet}{\parbox{\dimexpr\textwidth-2\fboxsep}{\strut\color{logoyellow}\rule[-0.35em]{0.8em}{1.35em}\hspace{0.6em}\color{white}\thesection.\ #1}}}
\titlespacing*{\section}{0pt}{1.6em}{0.8em}

% A pixel "star" like the ones in the menu background.
\newcommand{\pixstar}[3]{% x (mm), y (mm), pixel size (mm)
  \fill[stargrey] ({#1*1mm},{#2*1mm}) rectangle ++({#3*1mm},{#3*1mm});
  \fill[stargrey] ({(#1-#3)*1mm},{(#2-#3)*1mm}) rectangle ++({#3*1mm},{#3*1mm});
  \fill[stargrey] ({(#1+#3)*1mm},{(#2-#3)*1mm}) rectangle ++({#3*1mm},{#3*1mm});
  \fill[stargrey] ({#1*1mm},{(#2-2*#3)*1mm}) rectangle ++({#3*1mm},{#3*1mm});}
\newcommand{\key}[1]{\texttt{#1}}
\newcommand{\hd}[1]{\textcolor{white}{\textbf{#1}}}   % table header cell (navy row)
\newcolumntype{L}{>{\raggedright\arraybackslash}X}
\setlength{\parindent}{0pt}
\setlength{\parskip}{0.6em}
\fvset{frame=leftline,rulecolor=\color{violet},framerule=1mm,framesep=3mm,fontsize=\small}

\pagestyle{fancy}
\fancyhf{}
\fancyhead[L]{\small\bfseries\color{violet}for\_astronaut}
\fancyhead[R]{\small for\_code UFRJ}
\fancyfoot[C]{\small\bfseries\color{navy}\thepage}

\begin{document}

% ------------------------------------------------------------------ cover
% Same look as the game's title screen: navy background, pixel stars,
% white logo with a yellow "_" and the pixel-art astronaut.
\begin{titlepage}
\begin{tikzpicture}[remember picture, overlay, shift={(current page.south west)}]
  \fill[navy] (0,0) rectangle (210mm,297mm);
  \pixstar{18}{275}{6}  \pixstar{150}{267}{4} \pixstar{172}{227}{6}
  \pixstar{34}{179}{4}  \pixstar{120}{147}{6} \pixstar{182}{101}{4}
  \pixstar{12}{65}{6}   \pixstar{96}{35}{4}   \pixstar{60}{237}{4}
  \node[anchor=north west, text width=150mm, inner sep=0] at (22mm,215mm) {%
    {\color{logoyellow}\bfseries\sffamily\small U\,S\,E\,R\quad M\,A\,N\,U\,A\,L}\\[4mm]
    {\color{white}\bfseries\sffamily\fontsize{40}{44}\selectfont for{\color{logoyellow}\_}astronaut}\\[7mm]
    \colorbox{violet}{\parbox{110mm}{\vspace{2mm}\color{white}\sffamily
      How to install, start and play for\_astronaut, the space-mission design
      card game: plan a science mission, build the spacecraft, pick an orbit
      and a launcher, fly it and read your score.\vspace{2mm}}}};
  \node[anchor=south east, inner sep=0] at (188mm,40mm)
    {\includegraphics[width=52mm]{assets/images/astronaut/waving.png}};
  \node[anchor=south west, inner sep=0] at (22mm,16mm)
    {\color{softgrey}\sffamily\scriptsize for\_code UFRJ $\cdot$ NASA Space Apps Challenge 2026 $\cdot$ LaTeX version of \texttt{docs/manual.md}};
\end{tikzpicture}
\end{titlepage}

{\hypersetup{linkcolor=navy}\tableofcontents}
\clearpage

\section{Play online}

Open \url{https://buenosaas.itch.io/for-astronaut} in a browser and press
\textbf{Play}. Nothing needs to be installed.

\section{Run on your computer}

You need Python 3.10 or newer.

\begin{Verbatim}
python -m pip install -r requirements.txt
python for_astronaut.py
\end{Verbatim}

A 1000x760 window opens with the title screen and the theme music.

To try the landscape mobile layout on a computer:

\begin{Verbatim}
python for_astronaut_mobile/main.py
\end{Verbatim}

\section{Controls}

\begin{tabularx}{\textwidth}{@{}LLLL@{}}
\toprule
\rowcolor{navy}\hd{Action} & \hd{Mouse} & \hd{Keyboard} & \hd{Touch (mobile)} \\
\midrule
Start from the title screen & Click \textbf{PLAY} & \key{Space} or \key{Enter} & Tap \textbf{PLAY} \\
Open the ``How it works'' guide (desktop) & Click \textbf{HOW IT WORKS} & --- & --- \\
Close the guide & Click anywhere & \key{Space}, \key{Enter} or \key{Esc} & --- \\
Preview an option & Move the mouse to the left or right half & $\leftarrow$ / $\rightarrow$ or \key{A} / \key{D} & Tap a side \\
Choose the previewed option & Click & Press the same key again & Tap the same side again \\
Continue a briefing or a finished flight & Click & \key{Space} or \key{Enter} & Tap \\
Restart on the results screen & --- & \key{R} & --- \\
Quit & Close the window & \key{Esc} (on any screen except the guide) & --- \\
\bottomrule
\end{tabularx}

\section{How a mission works}

New to the game? On the title screen, \textbf{HOW IT WORKS} opens a one-page
guide to the trade-off at the heart of every mission: cost, mass and schedule.

Each card shows two options. Each option has a good \textbf{and} a bad side,
and the numbers on the card tell you how it changes your budget, mass, power,
schedule and science.

\begin{enumerate}[leftmargin=*]
  \item \textbf{Mission briefing}: choose the science objective and the program
        difficulty. The briefing sets your budget, dry-mass limit and deadline.
  \item \textbf{Spacecraft assembly}: pick the structure, power, attitude
        control, computer, communications and instruments. Watch the status
        bar: dry mass against its limit, cost against budget, power used
        against power generated.
  \item \textbf{Orbit and launcher}: choose the orbit (low Earth, polar,
        geostationary or lunar) and the launcher. The compatibility check
        compares what the launcher can lift to that orbit with your
        spacecraft's mass at liftoff.
        \begin{itemize}
          \item \textbf{PASSED}: you see the margin and how well the orbit fits
                your objective.
          \item \textbf{NOT COMPATIBLE}: the craft is too heavy for that
                launcher and orbit.
        \end{itemize}
  \item \textbf{Mission simulation}: watch the launch. Empty stages drop away
        and the rocket gets lighter. Then handle the events of the flight,
        including \textbf{critical events} (marked with an alarm) whose outcome
        is hidden until you choose. Keep an eye on battery, propellant and
        onboard data.
  \item \textbf{Results}: read your scores and the feedback on what worked and
        what did not.
\end{enumerate}

\section{Reading the results}

\begin{tabularx}{\textwidth}{@{}lL@{}}
\toprule
\rowcolor{navy}\hd{Score} & \hd{What raises it} \\
\midrule
Science (40\%)    & Instruments, an orbit that fits the objective, and data actually sent to Earth \\
Cost (20\%)       & Money left in the budget \\
Safety (25\%)     & Reliability and margins of power, mass, launcher capacity and fuel \\
Creativity (15\%) & Bold or unusual choices \\
\bottomrule
\end{tabularx}

The weighted total gives a grade: \textbf{S} (85 or more), \textbf{A} (70),
\textbf{B} (55), \textbf{C} (40) or \textbf{D} (below 40). A failed mission
halves the total. A spacecraft that never launches gets \textbf{F}.

\section{Tips}

\begin{itemize}[leftmargin=*]
  \item More mass needs more fuel, which needs a bigger rocket, which costs
        more and takes longer. Saving on one of cost, mass or schedule usually
        costs you on another.
  \item A powerful instrument is useless if the antenna cannot send its data
        home.
  \item Leave margin: a craft right at the mass or power limit has nothing left
        when an event goes wrong.
\end{itemize}

\section{Common problems}

\begin{tabularx}{\textwidth}{@{}LL@{}}
\toprule
\rowcolor{navy}\hd{Problem} & \hd{Fix} \\
\midrule
\texttt{ModuleNotFoundError: No module named 'pygame'} & Run \texttt{python -m pip install -r requirements.txt} \\
\texttt{FileNotFoundError} for a \texttt{.png} file at start-up & The \texttt{assets/} folder is missing or incomplete; download the whole repository \\
The game runs but is silent & The files in \texttt{assets/audio/} are missing, or the computer has no audio output; the game still works without sound \\
The window does not fit the screen & Use the online version or the mobile layout, which scales to the screen \\
\bottomrule
\end{tabularx}

\vfill
\noindent\rule{\textwidth}{0.4pt}

{\small\itshape This manual was written with AI assistance (Claude Code) and
reviewed by the team. The PDF version is generated from
\texttt{docs/manual.md}; this LaTeX version has the same content. See
\texttt{docs/ai\_usage.md}.}

\end{document}
